\textbf{Phần 1. Trắc nghiệm lựa chọn}
\setcounterpageref{ex}{0}

\begin{ex} %[4D13N67]
Số \textit{Một trăm hai mươi lăm triệu} viết là
\choice
{$12\,500\,000$}
{\True $125\,000\,000$}
{$1\,250\,000$}
{$125\,000$}
\loigiai{
Số đó là $125\,000\,000$.\\
\textbf{Đáp án của câu hỏi là B}
}
\end{ex}

\begin{ex} %[4D13N67]
Trong số $45\,678\,900$, chữ số 5 ở hàng nào?
\choice
{Hàng chục triệu}
{Hàng trăm nghìn}
{\True Hàng triệu}
{Hàng chục nghìn}
\loigiai{
Chữ số 5 ở hàng triệu.\\
\textbf{Đáp án của câu hỏi là C}
}
\end{ex}

\begin{ex} %[4D13H67]
Làm tròn số $3\,480\,000$ đến hàng trăm nghìn ta được
\choice
{\True $3\,500\,000$}
{$3\,400\,000$}
{$3\,000\,000$}
{$4\,000\,000$}
\loigiai{
Hàng chục nghìn là 8 nên làm tròn lên $3\,500\,000$.\\
\textbf{Đáp án của câu hỏi là A}
}
\end{ex}

\begin{ex} %[4D13H67]
Số liền sau của số $99\,999\,999$ là
\choice
{$99\,999\,998$}
{$100\,000\,001$}
{$10\,000\,000$}
{\True $100\,000\,000$}
\loigiai{
$99\,999\,999 + 1 = 100\,000\,000$.\\
\textbf{Đáp án của câu hỏi là D}
}
\end{ex}

\begin{ex} %[4D13H67]
Số lớn nhất có 8 chữ số khác nhau là
\choice
{$99\,999\,999$}
{$98\,765\,430$}
{\True $98\,765\,432$}
{$87\,654\,321$}
\loigiai{
Số lớn nhất có 8 chữ số khác nhau là $98\,765\,432$.\\
\textbf{Đáp án của câu hỏi là C}
}
\end{ex}

\begin{ex} %[4D13H67]
Số nào sau đây chia hết cho cả 2 và 5?
\choice
{$1235$}
{$2468$}
{$3579$}
{\True $4560$}
\loigiai{
Số tận cùng bằng 0 chia hết cho cả 2 và 5.\\
\textbf{Đáp án của câu hỏi là D}
}
\end{ex}

\begin{ex} %[4D13V67]
Sắp xếp các số sau theo thứ tự từ lớn đến bé: $45\,020\,000; 45\,200\,000; 45\,002\,000; 45\,220\,000$.
\choice
{\True $45\,220\,000 > 45\,200\,000 > 45\,020\,000 > 45\,002\,000$}
{$45\,002\,000 > 45\,020\,000 > 45\,200\,000 > 45\,220\,000$}
{$45\,200\,000 > 45\,220\,000 > 45\,020\,000 > 45\,002\,000$}
{$45\,220\,000 > 45\,020\,000 > 45\,200\,000 > 45\,002\,000$}
\loigiai{
Thứ tự: $45\,220\,000 > 45\,200\,000 > 45\,020\,000 > 45\,002\,000$.\\
\textbf{Đáp án của câu hỏi là A}
}
\end{ex}

\begin{ex} %[4D13V67]
Tìm chữ số $x$ để số $\overline{35x,420} < 352\,420$:
\choice
{$x = 2$}
{\True $x \in \{0; 1\}$}
{$x = 3$}
{$x = 4$}
\loigiai{
$x < 2 \Rightarrow x \in \{0; 1\}$.\\
\textbf{Đáp án của câu hỏi là B}
}
\end{ex}

\begin{ex} %[4D13V67]
Dãy số tự nhiên chẵn liên tiếp bắt đầu từ 2 đến 100 có bao nhiêu số hạng?
\choice
{$49$ số}
{$51$ số}
{$100$ số}
{\True $50$ số}
\loigiai{
$(100 - 2) : 2 + 1 = 50$ số.\\
\textbf{Đáp án của câu hỏi là D}
}
\end{ex}

\begin{ex} %[4D13V67]
Tính tổng của dãy số: $1 + 2 + 3 + \dots + 30$.
\choice
{$450$}
{$480$}
{\True $465$}
{$435$}
\loigiai{
$(1 + 30) \times 30 : 2 = 465$.\\
\textbf{Đáp án của câu hỏi là C}
}
\end{ex}

\begin{ex} %[4D13V67]
Trong số $789\,123\,456$, tổng các chữ số thuộc lớp triệu là
\choice
{$18$}
{\True $24$}
{$15$}
{$21$}
\loigiai{
Lớp triệu gồm $7, 8, 9$. Tổng: $7 + 8 + 9 = 24$.\\
\textbf{Đáp án của câu hỏi là B}
}
\end{ex}

\begin{ex} %[4D13C67]
Có bao nhiêu số có 6 chữ số khác nhau được viết từ các chữ số $0, 1, 2, 3, 4, 5$?
\choice
{\True $600$ số}
{$720$ số}
{$120$ số}
{$500$ số}
\loigiai{
Chữ số đầu tiên có 5 cách chọn (khác 0). Các vị trí tiếp theo có $5 \times 4 \times 3 \times 2 \times 1 = 120$ cách. Tổng cộng: $5 \times 120 = 600$ số.\\
\textbf{Đáp án của câu hỏi là A}
}
\end{ex}

\textbf{Phần 2. Trắc nghiệm đúng sai}
\setcounterpageref{ex}{0}

\begin{ex} %[4D13H67]
Xét tính đúng sai về số tự nhiên:
\choiceTF
{\True Số 0 là số tự nhiên bé nhất.}
{\True Không có số tự nhiên lớn nhất.}
{Hai số tự nhiên liên tiếp hơn kém nhau 2 đơn vị.}
{\True Số liền trước của 1 triệu là $999\,999$.}
\loigiai{
\begin{itemize}
    \item 0 là số tự nhiên bé nhất: Đúng.
    \item Không có số lớn nhất: Đúng.
    \item Hai số tự nhiên liên tiếp hơn kém nhau 1 đơn vị (không phải 2). Sai.
    \item $1\,000\,000 - 1 = 999\,999$: Đúng.
\end{itemize}
\begin{itemize}
    \item a) \textbf{Đúng}.
    \item b) \textbf{Đúng}.
    \item c) \textbf{Sai}.
    \item d) \textbf{Đúng}.
\end{itemize}
}
\end{ex}

\begin{ex} %[4D13V67]
Cho số $A = 50\,000\,000 + 700\,000 + 4000 + 8$. Xét tính đúng sai:
\choiceTF
{\True Số $A$ được viết là $50\,704\,008$.}
{\True Lớp triệu của số $A$ là $50$.}
{Chữ số 7 có giá trị là $70\,000$.}
{\True Chữ số 7 có giá trị là $700\,000$.}
\loigiai{
\begin{itemize}
    \item Viết số đúng: $50\,704\,008$.
    \item Lớp triệu: 50. Đúng.
    \item Chữ số 7 ở hàng trăm nghìn nên có giá trị $700\,000 \neq 70\,000$. Ý 3 Sai, Ý 4 Đúng.
\end{itemize}
\begin{itemize}
    \item a) \textbf{Đúng}.
    \item b) \textbf{Đúng}.
    \item c) \textbf{Sai}.
    \item d) \textbf{Đúng}.
\end{itemize}
}
\end{ex}

\textbf{Phần 3. Trả lời ngắn}
\setcounterpageref{ex}{0}

\begin{ex} %[4D13H67]
Viết số: Tám mươi tư triệu ba trăm nghìn.
\shortans{84300000}
\loigiai{
\textbf{Đáp số:} $84300000$\\
$84\,300\,000$.
}
\end{ex}

\begin{ex} %[4D13H67]
Làm tròn số $7\,650\,000$ đến hàng trăm nghìn.
\shortans{7700000}
\loigiai{
\textbf{Đáp số:} $7700000$\\
$7\,700\,000$.
}
\end{ex}

\begin{ex} %[4D13V67]
Tìm số chẵn lớn nhất có 7 chữ số khác nhau.
\shortans{9876542}
\loigiai{
\textbf{Đáp số:} $9876542$\\
$9\,876\,542$.
}
\end{ex}

\begin{ex} %[4D13V67]
Trong số $456\,789\,123$, tính tổng các chữ số thuộc lớp nghìn.
\shortans{24}
\loigiai{
\textbf{Đáp số:} $24$\\
$7 + 8 + 9 = 24$.
}
\end{ex}

\begin{ex} %[4D13V67]
Dãy số $5, 10, 15, \dots, 100$ có bao nhiêu số hạng?
\shortans{20}
\loigiai{
\textbf{Đáp số:} $20$\\
$(100 - 5) : 5 + 1 = 20$.
}
\end{ex}

\begin{ex} %[4D13C67]
Tính tổng: $1 + 2 + 3 + \dots + 100$.
\shortans{5050}
\loigiai{
\textbf{Đáp số:} $5050$\\
$(1 + 100) \times 100 : 2 = 5050$.
}
\end{ex}
